\chapter[固有态射的有限性与仿射上同调维数]{固有态射的有限性定理；仿射代数概形的上同调维数}
%\addtocontents{toc}{M. Artin}
\begin{center}
M. Artin
\end{center}
\phantomsection\MakeLinkTarget*{sga4.exp.XIV}\label{sga4.exp.XIV}

\setcounter{pageoriginal}{144}

\newcommand{\degtr}{{\rm deg.tr.\ }}

本\pageoriginale 讲包含两个重要定理，其证明用到固有态射的基变换定理。

\section{固有态射的有限性定理}\etiquette[1]{XIV.sec1}

\begin{enonce*}{定理 1.1}\etiquette[1.1]{XIV.thm1.1}
设 $f:X\rightarrow S$ 是有限表现的固有态射（proper morphism）。设 $A$ 是左 Noether 环，并且对于某个适当的 $n\in\NN$ 是一个 $\ZZ/n$-代数。设 $F$ 是 $X$ 上的\emph{可构造}（constructible）集合层（\resp 群层，\resp $A$-模层）。那么层 $f_*F$（\resp $R^1f_*F$，\resp 对每个 $q\geq0$ 的 $R^qf_*F$）也是可构造的。
\end{enonce*}

特别地，我们有：

\begin{enonce*}{推论 1.2}\etiquette[1.2]{XIV.cor1.2}
设 $X$ 是可分闭域 $k$ 的谱上的固有概形，$F$ 是 $X$ 上的可构造挠阿贝尔层。那么对每个 $q\geq0$，群 $H^q(X,F)$ 都是有限群。
\end{enonce*}

\begin{enonce*}[remark]{注 1.3}\etiquette[1.3]{XIV.rem1.3}
如果去掉 $f$ 固有这一条件，定理对于阿贝尔层便不成立：取 $F=\ZZ/p$，$p=\car k$，$X=\EE^1_k$ 为仿射直线，并取 $f$ 为 $X$ 到 $\PP^1_k$ 的包含态射或结构态射 $X\rightarrow\Spec k$，即可看出这一点。然而，有如下猜想\nde{Ofer Gabber 在 2005 年纪念 Luc Illusie 和 Pierre Deligne 的会议上报告了 \hyperref[XIV.thm1.1]{1.1} 的以下推广。设 $f:X\to Y$ 是拟优 Noether 概形之间的有限型态射（\EGA IV 7.8），$F$ 是 $X_{\text{ét}}$ 上的可构造 $\mathbf Z/n$-模层，其中 $n$ 在 $Y$ 上可逆。那么 $Y_{\text{ét}}$ 上的层 $R^if_*F$ 都是可构造的，并且除有限多个 $i$ 外均为零。对于可构造集合层以及阶在 $Y$ 上可逆的群层，也有类似的断言。\newline
在 Gabber 的结果之前，\hyperref[XIV.thm1.1]{定理 1.1} 已在以下情形中得到证明。设 $S$ 是维数 $\leq1$ 的正则概形，$f:X\to Y$ 是 $S$-概形之间的有限型态射，$F$ 是 $X_{\text{ét}}$ 上的左 $A$-模层。那么层 $R^if_*F$ 都是可构造的（\SGA 4 1/2, Finitude 1.1）。}：若 $F$ 的阶与 $S$ 的\pageoriginale 剩余特征互素，且 $S$ 是“优的”（excellent）（\EGA IV 7.8），则该定理成立。当奇点消解可用时（例如特征为零时），对于优的等特征概形 $S$ 可以证明这一点（\hyperrefext[XIX][XIX.thm5.1]{5.1}）。
\end{enonce*}

\noindent
\hyperref[XIV.thm1.1]{1.1} 的\emph{证明}。断言在 $S$ 上是局部的，所以可取 $S$ 为仿射概形。于是 $S$ 是 $\ZZ$ 上有限型环的谱的极限，而定理中的资料可以由一个固有态射 $f_0:X_0\rightarrow S_0$ 以及 $X_0$ 上的一个可构造层 $F_0$，经基变换 $S\rightarrow S_0$ 得到，其中 $S_0$ 在 $\ZZ$ 上有限型；见 \EGA IV 8 和 \hyperrefext[IX][IX.cor2.7.4]{2.7.4}。由于可构造层的逆像仍然可构造（\hyperrefext[IX][IX.prop2.4]{2.4} (iii)），将 \hyperrefext[XII][XII.thm5.1]{5.1} 应用于基变换态射 $S_0\leftarrow S$，便将问题归结为证明 $f_{0*}F_0$（\resp $R^qf_{0*}F_0$）的可构造性，由此得到下面的引理。

\begin{enonce*}{引理 1.4}\etiquette[1.4]{XIV.lem1.4}
只需在 $S$ 是 $\Spec\ZZ$ 上的有限型概形时验证 \hyperref[XIV.thm1.1]{1.1}。
\end{enonce*}

\begin{enonce*}{引理 1.5}\etiquette[1.5]{XIV.lem1.5}
\hyperref[XIV.thm1.1]{1.1} 中关于集合的断言成立：即若 $F$ 是可构造集合层，则 $f_*F$ 也是。
\end{enonce*}

\noindent
\emph{证明}。可设 $S$ 是 Noether 概形（\hyperref[XIV.lem1.4]{1.4}）。取一个单射 $F\rightarrow G=\prod_i\pi_{i*}(C_i)$，其中 $\pi_i:X_i\rightarrow X$ 是有限态射，且 $C_i$ 在 $X_i$ 上常值且可构造（\hyperrefext[IX][IX.prop2.14]{2.14}）。有 $f_*F\subset f_*G$；因此，由 \hyperrefext[IX][IX.prop2.9]{2.9} (ii)，只需处理 $F=G$ 的情形。现在 $f_*G=\prod_i(f\pi_i)_*C_i$，而 $S$ 上可构造层的有限积仍然可构造（\hyperrefext[IX][IX.prop2.6]{2.6} (ibis)）。用 $X_i$ 代替 $X$，便归结为 $F=D_X$ 是取值为 $D$ 的可构造常值层的情形。由 \hyperrefext[XII][XII.thm5.1]{5.1} (i)，纤维 $(f_*F)_{\bar s}$ 同构于 $D^{C(s)}$，其中 $C(s)=\pi_0(X_{\bar s})$；\pageoriginale 因而它是有限集。由 \hyperrefext[IX][IX.prop2.13]{2.13} (iii)，只需证明函数“$f_*F$ 的纤维中的元素个数”在 $S$ 上是可构造函数。这个函数为 $\card(D^C)$，其中 $C$ 是函数“几何纤维中的连通分支个数”，而 $C$ 是可构造函数（\EGA IV 9.7.9），这就证明了结论。

\begin{enonce*}{引理 1.6}\etiquette[1.6]{XIV.lem1.6}
设 $S$ 是 Noether 概形，$f:X\rightarrow S$ 固有，$F$ 是 $X$ 上的可构造群层（\resp $A$-模层），$q$ 为等于 $0$ 或 $1$ 的整数（\resp $\geq0$ 的整数）。设 $S_i$ $(1\leq i\leq n)$ 为 $S$ 的子概形，满足 $S=\bigcup_iS_i$；设 $F_i$ 和 $f_i:X_i\rightarrow S_i$ 是这些资料在 $S_i$ 上的“限制”。若对于每个 $i$，$R^qf_{i*}F_i$ 都可构造，则 $R^qf_*F$ 也可构造。
\end{enonce*}

\noindent
\emph{证明}。由 \hyperrefext[XII][XII.thm5.1]{5.1}，$R^qf_*F$ 在 $S_i$ 上的限制同构于 $R^qf_{i*}F_i$，而引理由 \hyperrefext[IX][IX.prop2.8]{2.8} 得出。

现在采用与 \hyperrefext[XII][XII.sec8]{8} 类似的论证：

\begin{enonce*}{引理 1.7}\etiquette[1.7]{XIV.lem1.7}
设 $S$ 是 Noether 概形，
$$
\xymatrix@=1.25cm{
X        \ar[rd]_{f} &           &\bar{X} \ar[ll]_{\pi}
 \ar[ld]^{\bar{f}}\\
                 &    S      &
}
$$
是固有态射的交换图，$j:U\rightarrow X$ 是一个开子概形，使得开子概形 $\bar U=\fprod{U}{\bar X}{X}$ 同构地映到 $U$。设 $Y=X-U$ 是约化闭子概形，$f_0:Y\rightarrow S$ 和 $i:Y\rightarrow X$ 是典范态射。若\hyperref[XIV.thm1.1]{定理 1.1} 对 $\bar f$、$f_0$ 和 $\pi$ 成立\nde{关于 $R^q\pi_*$ 的假设仅在 $q=0$ 时用于证明。在这一情形，可构造性已由 \hyperref[XIV.lem1.5]{1.5} 得知。}，则它对 $f$ 也成立。
\end{enonce*}

\noindent
\emph{证明}。\pageoriginale 先考虑 $X$ 上的 $A$-模层 $F$。有正合列
$$
0\longrightarrow j_!j^*F\longrightarrow F\longrightarrow
i_*i^*F\longrightarrow0{\quoi}.
$$
由相应的 $R^qf_*$ 的上同调正合列以及 \hyperrefext[IX][IX.prop2.6]{2.6}，只需验证两端项的可构造性。注意，由 \hyperrefext[IX][IX.prop2.4]{2.4} (iii) 和 \hyperrefext[IX][IX.prop2.14]{2.14} (i)，$i_*i^*F$ 可构造，因而 $j_!j^*F$ 也可构造（\hyperrefext[IX][IX.prop2.6]{2.6}）。由于 $i$ 是闭浸入，由 \hyperrefext[VIII][VIII.cor5.6]{5.6}，
$$
(R^qf_*)i_*(i^*F)\simeq(R^qf_{0*})(i^*F){\quoi},
$$
因而这一情形的可构造性由对 $f_0$ 的假设得出。此外，若 $\bar j:\bar U\rightarrow\bar X$ 表示包含态射，且 $\bar F=\pi^*F$，则
\begin{equation*}
j_!j^*F\simeq\pi_*\bar j_!\bar j^*\bar F{\quoi},\tag{$*$}
\end{equation*}
为此，先定义一个从左到右的态射，再利用 \hyperrefext[XII][XII.thm5.1]{5.1} 逐纤维证明其可逆即可。在 $Y$ 的每个几何点 $\bar y$ 处，$\bar j_!\bar j^*\bar F$ 在 $\bar X$ 关于 $X$ 的纤维 $\bar X_{\bar y}$ 上的限制为零；因此，由 \hyperrefext[XII][XII.thm5.1]{5.1} (iii)，对于每个 $q\geq0$，$(R^q\pi_*)(\bar j_!\bar j^*\bar F)$ 在 $\bar y$ 处的纤维为零。另一方面，$\pi$ 在 $X-Y$ 上诱导一个同构，所以对于 $q>0$，
$$
(R^q\pi_*)(\bar j_!\bar j^*\bar F)=0
$$
。因此，Leray 谱序列
$$
E^{p,q}_2=(R^pf_*)(R^q\pi_*)(\bar j_!\bar j^*\bar F)
\Longrightarrow R^n\bar f_*(\bar j_!\bar j^*\bar F)
$$
与 $(*)$ 一起给出同构
$$
(R^q\bar f_*)(\bar j_!\bar j^*\bar F)\simeq
(R^qf_*)(j_!j^*F)
$$
，对所有 $q\geq0$ 成立。由于 $\bar j_!\bar j^*\bar F$ 可构造（与 $j_!j^*F$ 的论证相同），由对 $\bar f$ 的假设，左端可构造。于是 $(R^qf_*)(j_!j^*F)$ 可构造，因而 $R^qf_*F$ 也可构造。\pageoriginale

现在设 $F$ 是 $X$ 上的可构造群层。用 $\pi\amalg i:(\bar X\amalg Y)\rightarrow X$ 代替 $\pi:\bar X\rightarrow X$。于是只需证明：若 $\pi:\bar X\rightarrow X$ 是\emph{满射}，且 \hyperref[XIV.thm1.1]{1.1} 对 $\pi$ 和 $\bar f$ 成立，则它对 $f$ 也成立。先设 $\bar F$ 是 $\bar X$ 上的可构造群层。那么 $\pi_*\bar F$ 在 $X$ 上可构造（\hyperref[XIV.lem1.5]{1.5}），且有单射
$$
(R^1f_*)\pi_*\bar F\hookrightarrow(R^1\bar f_*)\bar F.
$$
由假设右端可构造，因而左端也可构造（\hyperrefext[IX][IX.prop2.9]{2.9} (ii)）。由于 $\pi$ 满射，有单射 $F\rightarrow\pi_*\pi^*F=G$。层 $\bar F=\pi^*F$ 可构造，所以由刚才的证明，$G$ 和 $R^1f_*G$ 也可构造。这样就归结为下面的引理。

\begin{enonce*}{引理 1.8}\etiquette[1.8]{XIV.lem1.8}
设 $S$ 是 Noether 概形，$f:X\rightarrow S$ 是固有态射，而 $\displaystyle F\mathop{\hookrightarrow}^{u}G$ 是 $X$ 上可构造群层的单射。若 $R^1f_*G$ 可构造，则 $R^1f_*F$ 也可构造。
\end{enonce*}

\noindent
\emph{证明}。令 $C=G/F$，这是一个可构造的 $G$-齐性空间层（应用 \hyperrefext[IX][IX.prop2.6]{2.6}）。考虑正合列
\begin{equation*}
f_*C\xrightarrow{\delta}R^1f_*F\xrightarrow{u^1}R^1f_*G{\quoi}.
\tag{$*$}
\end{equation*}
为证明 $R^1f_*F$ 可构造，根据 Noether 归纳法和 \hyperref[XIV.lem1.6]{1.6}，只需证明：若 $S\neq\varnothing$，则 $S$ 有一个非空开子集，使 $R^1f_*F$ 在其上可构造。因而可设 $R^1f_*G$ 局部常值。由于可构造性关于平展（étale）拓扑是局部性质，还可进一步假设 $S$ 不可约，并且 $R^1f_*G$ 是取值为 $E=\{e_1,\ldots,e_n\}$ 的常值（且可构造）层。对于 $i=1,\ldots,n$，设 $D_i$ 是 $R^1f_*F$ 的子层\pageoriginale，它是截面 $e_i$ 在 $u^1$ 下的逆像。由于拓扑斯（topos）中“余积是泛的”，有 $R^1f_*F=\amalg_iD_i$；因而只需证明每个 $D_i$ 在某个适当的非空 $S$-平展概形 $S'$ 上可构造。

回忆 $R^1f_*F$ 是预层 $\Rr^1F$ 的伴随层，其中 $\Rr^1F(S')=H^1(X\times_SS',F)$。因此，可以假设（用非空的 $S$-平展概形 $S'$ 代替 $S$ 之后），$D_i$ 是“空层”（从而可构造），\emph{或}存在一个元素 $\alpha\in H^{1}(X,F)$，诱导 $D_i$ 的截面 $\sisi{\alpha}{\bar{\alpha}}$。在后一情形，设
$$
f_{\ast}C^{\alpha}\xrightarrow{\delta}R^{1}f_{\ast}F^{\alpha}\longrightarrow
R^{1}f_{\ast}G^{\alpha}
$$
是对 $(\ast)$ 用 $\alpha$ 作扭曲所得到的正合列。有
$$
R^{1}f_{\ast}F^{\alpha}\simeq R^{1}f_{\ast}F\quav
R^{1}f_{\ast}G^{\alpha}\simeq R^{1}f_{\ast}G{\quoi},
$$
并且 $R^{1}f_{\ast}F^{\alpha}$ 中与 $\bar{\alpha}$ 对应的截面是单位截面。因此，$R^{1}f_{\ast}F^{\alpha}$ 中与 $D_i$ 对应的子层 $D_i^{\alpha}$ 是 $R^{1}f_{\ast}G^{\alpha}$ 的单位截面的逆像；即 $D_i^{\alpha}$ 是 $f_{\ast}C^{\alpha}$ 在 $\delta$ 下的像。由于 $f_{\ast}C^{\alpha}$ 可构造（\hyperref[XIV.lem1.5]{1.5}），$D_i^{\alpha}$ 也可构造，因而 $D_i$ 也可构造（应用 \hyperrefext[IX][IX.prop2.6]{2.6}）。这就证明了引理。

\begin{enonce*}{引理 1.9}\etiquette[1.9]{XIV.lem1.9}
设 $S$ 是 Noether 概形，
$$
\xymatrix@=1.25cm{
X       \ar[rd]_{f}      &      &\bar{X} \ar[ll]_{\pi} \ar[ld]^{\bar{f}}\\
              & S    &
}
$$
是固有态射的交换图。若定理 \hyperref[XIV.thm1.1]{1.1} 对\pageoriginale $f$ 和 $\pi$ 成立，则它对 $\bar{f}$ 也成立。
\end{enonce*}

\noindent
\emph{证明}。对于可构造 $A$-模层 $\bar{F}$，使用 Leray 谱序列
$$
E^{p,q}_2=(R^{p}f_{\ast})(R^{q}\pi_{\ast})\bar{F}\Longrightarrow
(R^{n}\bar{f}_{\ast})\bar{F}{\quoi}.
$$
由假设，对每个 $p,q$，$E^{p,q}_2$ 都可构造；于是由 \hyperrefext[IX][IX.prop2.6]{2.6}，收敛项也可构造。

设 $\bar{F}$ 是可构造群层。有正合列（\hyperrefext[XII][XII.prop3.2]{3.2}）
\begin{equation*}
0\longrightarrow (R^{1}f_{\ast})\pi_{\ast}\bar{F}\longrightarrow
(R^{1}\bar{f}_{\ast})\bar{F}\longrightarrow
f_{\ast}(R^{1}\pi_{\ast})\bar{F}{\quoi},\tag{$\ast$}
\end{equation*}
由假设和 \hyperref[XIV.lem1.5]{1.5}，两端项可构造。仿照 \hyperref[XIV.lem1.8]{1.8} 的证明：由 Noether 归纳法和 \hyperref[XIV.lem1.6]{1.6}，只需假设 $S\neq ∅$，并证明在某个适当的非空开子集上的可构造性。因此可设 $f_{\ast}(R^{1}\pi_{\ast})\bar{F}$ 局部常值，并且在作平展局部化之后，甚至是取值为 $E=\{e_1,\ldots,e_n\}$ 的常值可构造层（\sisi{\hyperrefext[IX][IX.prop2.8]{2.8}}{\hyperrefext[IX][IX.prop2.4]{2.4}~(i)}）。设 $D_i$ 是 $(R^{1}\bar{f}_{\ast})\bar{F}$ 中 $e_i$ 的逆像子层，于是 $(R^{1}\bar{f}_{\ast})\bar{F}=\coprod_iD_i$。只需证明每个层 $D_i$ 在某个适当的非空连通平展概形 $S'/S$ 上诱导可构造层。若 $D_i$ 不是“空层”，则可归结为存在元素 $\alpha\in H^{1}(\bar{X},\bar{F})$ 诱导 $D_i$ 的截面 $\bar{\alpha}$ 的情形。设
$$
\sisi{
0\longrightarrow (R^{1}f_{\ast})\pi_{\ast}\bar{F}\longrightarrow
(R^{1}\bar{f}_{\ast})\bar{F}\longrightarrow
f_{\ast}(R^{1}\pi_{\ast})\bar{F}
}
{0\longrightarrow (R^{1}f_{\ast})\pi_{\ast}\bar{F}^{\alpha}\longrightarrow
(R^{1}\bar{f}_{\ast})\bar{F}^{\alpha}\longrightarrow
f_{\ast}(R^{1}\pi_{\ast})\bar{F}^{\alpha}}
$$
是对 $(\ast)$ 中的 $\bar{F}$ 用 $\alpha$ 作扭曲得到的正合列。$(R^{1}\bar{f}_{\ast})\bar{F}^{\alpha}$ 中与 $\bar{\alpha}$ 对应的截面是单位截面。因此，\hbox{$(R^{1}\bar{f}_{\ast})\bar{F}^{\alpha}$} 中与 $D_i$ 对应的子层 $D_i^{\alpha}$ 是 $f_{\ast}(R^{1}\pi_{\ast})\bar{F}^{\alpha}$ 的单位截面的逆像，\pageoriginale 从而 $D_i^{\alpha}\simeq (R^{1}f_{\ast})\pi_{\ast}\bar{F}^{\alpha}$，这确实是一个可构造层，\cqfd

\begin{enonce*}{引理 1.10}\etiquette[1.10]{XIV.lem1.10}
只需在 $S$ 是 $\Spec\ZZ$ 上的有限型概形，且 $f$ 是相对维数 $\leq 1$ 的射影态射时验证 \hyperref[XIV.thm1.1]{1.1}。
\end{enonce*}

\noindent
\emph{证明}。先设 $f$ 是射影态射，并对 $f$ 的相对维数作归纳；用 $S$ 的一个仿射开子集代替它后，可设这一相对维数有限。若 $f$ 的相对维数 $\leq n$，其中 $n\geq 2$，那么局部地在 $S$ 上可以找到一个有限态射 $\pi:X\rightarrow\PP^n_S$。由 \hyperref[XIV.lem1.9]{1.9} 和 \sisi{IX 2.15}{\hyperrefext[IX][IX.prop2.14]{2.14}} (i)，归结为 $f$ 是 $\PP^n_S$ 的结构态射的情形。考虑态射的交换图
$$
\xymatrix@C=2cm{
                   & \bar{X}\ar[ld]_{\sisi{}{\pi}} \ar[rd]^{a} \ar@{-}[dd]^{\bar{f}} &       \\
\PP^{n}_S \ar[rd]_{f}  &                  &\PP^{1}_S \ar[ld]^{b}\\
                   &  S                          &
}
$$
，它在 \hyperrefext[XII][XII.lem8.3]{8.3} 的证明中已经考虑过，其中 $\bar{X}$ 由 $\PP^n_S$ 沿闭子概形 $Y\simeq\PP^{n-2}_S$ 作爆破得到。$b$ 的相对维数 $\leq1$，$a$ 的相对维数 $\leq n-1$；所以由归纳假设，定理对 $a$ 和 $b$ 成立，因而由 \hyperref[XIV.lem1.9]{1.9}，也对 $\bar{f}$ 成立。此外，$\pi$ 的相对维数 $\leq1$，而 $f_0:Y\rightarrow S$ 的相对维数 $\leq n-1$。因此由 \hyperref[XIV.lem1.7]{1.7}，\hyperref[XIV.thm1.1]{1.1} 对 $f$ 成立，从而对每个射影态射都成立。现在令 $f$ 为任意固有态射，对 $X$ 作 Noether 归纳：可设定理对 $X$ 的每个真闭子概形都成立，并且 $X\neq∅$。设
$$
\xymatrix@=1.5cm{
X  \ar[rd]_{f}          &           &\bar{X} \ar[ll]_{\pi} \ar[ld]^{\bar{f}}\\
                        &    S       &
}
$$
是 Chow 引理 \hyperrefext[XII][XII.lem7.1]{7.1} 给出的那类\pageoriginale 图。对这个图应用 \hyperref[XIV.lem1.7]{1.7}，即得 $f$ 满足 \hyperref[XIV.thm1.1]{1.1}，引理得证。

\setcounter{subsection}{10}
\subsection{}\etiquette[1.11]{XIV.subsec1.11}
现在可以完成 \hyperref[XIV.thm1.1]{1.1} 的证明。若 $F$ 是可构造群层，可以找到一个单射 $F\hookrightarrow G=\prod_i\pi_{i\sisi{}{,}\ast}C_i$，其中每个 $\pi_i$ 有限，每个 $C_i$ 常值且可构造（\sisi{IX 2.15}{\hyperrefext[IX][IX.prop2.14]{2.14}}）；因此由 \hyperref[XIV.lem1.8]{1.8}，可设 $F$ 常值。同样，若 $F$ 是可构造 $A$-模层，则应用 \sisi{IX 2.15}{\hyperrefext[IX][IX.prop2.14]{2.14}} 可以找到 $F$ 的一个消解 $F\rightarrow G^{\bigdot}$，$G^{\bigdot}=\{0\rightarrow G^0\rightarrow G^1\rightarrow G^2\rightarrow\ldots\}$，使每个 $G^{\nu}$ 具有形式 $G^{\nu}=\prod_i\sisi{\eta}{\pi}_{i\sisi{}{,}\ast}C_i$。消解 $G^{\bigdot}$ 的谱序列
$$
E^{p,q}_2=H^{\sisi{q}{p}}((R^{q}f_{\ast})G^{\bigdot})\Longrightarrow
(R^{n}f_{\ast})F{\quoi},
$$
便将问题归结为 $F$ 是取值为 $M$ 的常值可构造层的情形。

回忆 $A$ 是一个 $\ZZ/n$-代数。设 $\ell$ 是 $n$ 的一个素因子。有正合列
$$
0\longrightarrow \ell M\longrightarrow M\longrightarrow M/\ell
M\longrightarrow 0{\quoi},
$$
而仍然只需对 $\ell M$ 和 $M/\ell M$ 证明 \hyperref[XIV.thm1.1]{1.1}。

一个简单的归纳便将问题归结为 $A$ 是 $\ZZ/\ell$-代数，$\ell\in\PP$，且 $M$ 是有限型 $A$-模的情形。于是 $M$ 是 $\ZZ/\ell$ 上的向量空间，而且对每个拟紧拟分离的 $Y$，典范态射
$$
H^q(Y,\ZZ/\ell)\otimes_{\ZZ/\ell}M\longrightarrow H^q(Y,M)
$$
都是双射（事实上，若 $M$ 的秩有限，这是显然的，而两端都与滤过正向极限交换\pageoriginale （\sisi{VI 1.6}{\hyperrefext[VI][VI.thm5.1]{5.1} 和 \hyperrefext[VII][VII.prop3.3]{3.3}}））。因此
$$
[R^qf_{\ast}(\ZZ/\ell)]\otimes_{\ZZ/\ell}M\xrightarrow{\sim}
R^qf_{\ast}(M){\quoi}.
$$
所以显然只需证明 $R^qf_{\ast}(\ZZ/\ell)$ 作为 $\ZZ/\ell$-模是可构造的；换言之，我们归结为 $A=\ZZ/\ell$ 的情形，以下均作此假设。

注意，对于任意 $F$，$f$ 的相对维数 $\leq0$（即 $f$ 有限）的情形，现在由 \hyperrefext[VIII][VIII.cor5.6]{5.6} 和 \hyperref[XIV.lem1.5]{1.5} 得出。

对 $S$ 作 Noether 归纳。由 \hyperref[XIV.lem1.6]{1.6}，可设 $S\neq∅$，并可用任意非空开子集代替 $S$。此外，可设 $S$ 整，泛点为 $s$。纤维 $X_s$ 是维数 $\leq1$ 的代数概形；因此存在 $k(s)=K$ 的纯不可分扩张 $K'$，使 $(X_s\otimes_KK')_{\text{\rm red}}$ 的正规化 $Y$ 在 $\Spec K'$ 上光滑\footnote[1]{\Cf \EGA IV 17.15.14.}。用非空开子集代替 $S$ 之后，可设存在有限、根性、满射的态射 $S'\rightarrow S$，使 $S'_s$ 与 $\Spec K'$ 为 $K$-同构；还存在光滑射影态射 $\bar{X}'\rightarrow S'$ 及有限满射的 $S'$-态射 $\pi:\bar{X}'\rightarrow X'=\fprod{X}{S'}{S}$，它在 $S$ 的泛点 $s$ 上的纤维上诱导 $Y\rightarrow X_s\otimes_KK'$（\EGA IV 9.6.1）。现在\sisi{}{，}由 \hyperrefext[VIII][VIII.thm1.1]{1.1}\sisi{}{，}作纯不可分扩张不影响问题。因此可用 $S'$ 代替 $S$，用 $X'$ 代替 $X$；也就是说，可设 $f$ 出现在图
$$
\xymatrix@=1.5cm{
X   \ar[rd]_{f}& &\bar{X} \ar[ll]_{\pi} \ar[ld]^{\bar{f}}\\
         &         S        &           ,
}
$$
中，其中\pageoriginale $\pi$ 有限且满射，$\bar f$ 光滑。此外，按构造，可设 $X$ 有一个闭子概形 $Y$，其泛纤维有限，且 $\pi$ 在 $X-Y$ 上诱导同构。再用一个非空开子集代替 $S$，便归结为 $f_0:Y\rightarrow S$ 也有限的情形。

由于 $F$ 常值，$\pi^{\ast}F$ 也常值，并且因 $\pi$ 满射，$F\rightarrow\pi_{\ast}\pi^{\ast}F$ 是单射。若 $F$ 是群层，只需证明 $(R^1f_{\ast})(\pi_{\ast}\pi^{\ast}F)\simeq (R^1\bar f_{\ast})(\pi^{\ast}F)$ 可构造（\hyperref[XIV.lem1.8]{1.8}）。若 $F$ 是 $\ZZ/\ell$-模层，有正合列
$$
0\longrightarrow F\longrightarrow\pi_{\ast}\pi^{\ast}F\longrightarrow
C\longrightarrow0{\quoi},
$$
其中 $C$ 的支集在 $Y$ 中，因为 $\pi$ 在 $X-Y$ 上诱导同构。由此，$R^qf_{\ast}C=R^qf_{0\sisi{}{,}\ast}C$ 可构造（且当 $q>0$ 时为零），因此仍然只需证明 $(R^qf_{\ast})(\pi_{\ast}\pi^{\ast}F)= (R^q\bar f_{\ast})(\pi^{\ast}F)$ 可构造。用 $\bar f$ 代替 $f$，用 $\pi^{\ast}F$ 代替 $F$，便归结为 $F$ 常值且可构造、$f$ 光滑的情形。

先处理常值群层 $\sisi{f}{F}=G_X$。$(R^1f_{\ast})G_X$ 的所有纤维都有限。事实上，为证明这一点，应用 \hyperrefext[XII][XII.cor5.2]{5.2} (ii)，可归结为 $S$ 是可分闭域的谱，且 $X$ 在 $S$ 上光滑、维数 $\leq1$ 的情形。在这一情形结论已知（\SGA 1 X 2.6）。因此只需证明特殊化态射是单射（\hyperrefext[IX][IX.prop2.13]{2.13} (ii)）。为此，归结为 $S$ 严格局部的情形；由基变换定理 \hyperrefext[XII][XII.thm5.1]{5.1} (ii)，甚至可用严格局部离散赋值环的谱代替 $S$。此外，可设 $X$ 连通。由于\pageoriginale $H^1(X,G)$ 分类以 $G$ 为群的 $X$ 的主 Galois 覆盖，只需验证：若 $X'\rightarrow X$ 是连通平展覆盖，则其几何纤维也连通。这很容易由 $f$ 光滑且固有、并具有几何连通纤维这一事实得出。

现在处理 $F=(\ZZ/\ell)_X$ 的情形。当 $q=0$ 和 $q=1$ 时，结论已知（\hyperref[XIV.lem1.5]{1.5} 及前面的论证）。\sisi{由 \hyperrefext[X][X.cor4.3]{4.3} 和}{由} \hyperrefext[XII][XII.corbis5.3]{5.3 bis}，只剩下 $q=2$ 的情形；而且，若 $S$ 的特征为 $\ell$，则 $R^2f_{\ast}(\ZZ/\ell)=0$。因此，用 $S$ 的一个适当非空开子集代替它，可归结为 $\ell$ 在 $S$ 上可逆的情形。于是，由 \hyperrefext[XII][XII.cor5.2]{5.2} \sisi{(ii)}{(iii)} 和 \sisi{X 4.7}{\hyperrefext[X][X.cor5.2]{5.2}}，$R^2f_{\ast}(\ZZ/\ell)$ 在 $s\in S$ 上方几何点 $\bar s$ 处的纤维同构于 $(\ZZ/\ell)^{c(s)}$，其中 $c(s)$ 为几何纤维 $X_{\bar s}$ 的一维不可约分支数目。函数 $c$ 在 $S$ 上可构造（\EGA IV 9.7.9），结论由此得出（\hyperrefext[IX][IX.prop2.13]{2.13} (iii)）。

\begin{enonce*}[remark]{注 1.11}\etiquette[1.11]{XIV.rem1.11}
若使用具有固有支集的上同调的形式体系，\hyperref[XIV.thm1.1]{1.1} 的证明会简单得多；这一形式体系将在 \hyperref[sga4.exp.XVII]{XVII} 中作为固有态射基变换定理 \hyperrefext[XII][XII.thm5.1]{5.1} 的直接推论加以建立。此外，\hyperref[XIV.thm1.1]{1.1} 的陈述可推广到如下情形：在 $\sisi{Y}{S}$ 上给定任意挠环层 $𝒜$，以及 $X$ 上的一个 $f^{\ast}(𝒜)$-模复形 $K^{\bigdot}$，满足某种“可构造性”条件；在 \hyperref[XIV.thm1.1]{1.1} 的特殊情形中，该条件仅表示 $K^{\bigdot}$ 的上同调层 $\Hh^i(K^{\bigdot})$ 是可构造 $A$-模，并在 $i$ 充分大时为零。
\end{enonce*}

\section{维数的一种变式}\etiquette[2]{XIV.sec2}\pageoriginale

\subsection{}\etiquette[2.1]{XIV.subsec2.1}
设 $X$ 是概形。对于 $x\in X$，记
\begin{equation*}
d(x)=\dim\bar{\{x\}}\tag{2.1.1}
\end{equation*}\etiquette[2.1.1]{XIV.eq2.1.1}
，其中 $\bar{\{x\}}$ 是 $\{x\}$ 的闭包。若 $F$ 是 $X$ 上的阿贝尔层，令
\begin{equation*}
d(F)=\sup \{d(x)\mid x\in X \quaet F_{\bar{x}}\neq 0\}{\quoi}.\tag{2.1.2}
\end{equation*}\etiquette[2.1.2]{XIV.eq2.1.2}
在局部情形中，需要对这一概念作如下修正：

\begin{enonce*}{定义 2.2}\etiquette[2.2]{XIV.def2.2}
设 $Y$ 是泛链状（universally catenary）Noether 局部环的谱（\EGA IV 5.6.2），$f:X\rightarrow Y$ 是有限型态射。设 $x\in X$，$y=f(x)$。令
$$
\delta(x)=\dim\bar{\{y\}}+\degtr k(x)/k(y){\quoi}.
$$
若 $F$ 是 $X$ 上的阿贝尔层，令
$$
\delta(F)=\sup \{\delta(x)\mid x\in X\quaet F_{\bar{x}}\neq 0\}{\quoi}.
$$
\end{enonce*}

\begin{enonce*}{命题 2.3}\etiquette[2.3]{XIV.prop2.3}
\begin{enumerate}
%%%
%%% The typesetting of this enumerate environment is terrible.
%%% The spacing between the dash and (i) is too big.
%%% Also, the items (ii), ... should better be without indent.
%%% If you use the package enum.sty, then you can build the necessary list environment.
%%%
\renewcommand{\theenumi}{\roman{enumi}}
\renewcommand{\labelenumi}{(\theenumi)}
\item 设 $i:X'\rightarrow X$ 是浸入，$x'\in X'$。那么 $\delta(i(x'))=\delta(x')$。
%%% new line break within formula in text: spend extra line to induce new line breaks?

\item 设 $X_{0}$ 是 $X/Y$ 的闭纤维。有
$$
\delta(x)=
\begin{cases}
d(x) & \mbox{若 } \bar{\{x\}}\cap X_{0}\neq ∅\\
d(x)+1 & \mbox{若 } \bar{\{x\}}\cap X_{0}=∅{\quoi};
\end{cases}
$$

\item 若 $f$ 固有，则 $\delta(x)=d(x)$。
\end{enumerate}
\end{enonce*}

\noindent
\emph{证明}。\pageoriginale 断言 (i) 直接由定义得出。由于当 $f$ 固有时 $\bar{\{x\}}\cap X_{0}$ 非空，(iii) 由 (ii) 得出。对于 (ii)，可用 $\bar{\{x\}}$ 代替 $X$，用 $\bar{\{y\}}$ 代替 $Y$（赋予它们诱导的约化结构），(i) 保证可以这样做。设 $x'\in X$，$y'=f(x')$。有公式（\EGA IV 5.6.1.1）
$$
\dim \Oo_{X,x'}=\dim
\Oo_{Y,y'}+\degtr k(x)/k(y) - \degtr k(x')/k(y'){\quoi}.
$$

设 $X_{0}\neq ∅$。显然，$\displaystyle{\mathop{\sup}_{x'}}$ $\{\dim \Oo_{X,x'}\}=\dim X$ 在取 $x'$ 为 $X_{0}$ 的一个闭点时达到，因此
$$
\dim X=\dim \Oo_{X,x'}=\dim Y + \degtr k(x)/k(y)=\delta(x){\quoi}.
$$
若 $X_{0}=∅$，显然上确界在取 $x'$ 为其纤维 $f^{-1}(y')$ 中的闭点、且 $\dim\bar{\{y'\}}=1$ 时达到。这样的点存在，因为满足 $\dim\bar{\{y'\}}=1$ 的 $y'\in Y$ 构成一个极稠密集合（\EGA IV 10.1.3，\sisi{10.5.9, 1.8.4}{10.3.1, 10.4.1, 10.5.9}）。于是
\begin{align*}
\dim X=\dim \Oo_{X,x'} &= \dim \Oo_{Y,y'}+\degtr k(x)/k(y)\\
&= \dim Y-1 +\degtr k(x)/k(y)=\delta(x) \sisi{=}{-} 1{\quoi},
\end{align*}
结论得证。

\begin{enonce*}{命题 2.4}\etiquette[2.4]{XIV.prop2.4}
设 $f:X\rightarrow Y$ 是域 $k$ 的谱上有限型概形之间的态射。设 $x\in X$，$y=f(x)$，$y_{0}$ 为 $y$ 的一个特殊化。设 $f':X'\rightarrow Y'$ 是 $f$ 在 $y_{0}$ 上方几何点 $\bar{y}_{0}$ 处的严格局部化，$x'\in X'$ 是 $x$ 上方的一个点。那么
$$
\delta(x')=d(x)-d(y_{0}){\quoi}.
$$
\end{enonce*}

\noindent
\emph{证明}。\pageoriginale 有
\begin{align*}
d(x) &= d(y)+\degtr k(x)/k(y)\\
\delta(x') &= d(y')+\degtr k(x')/k(y'){\quoi},
\end{align*}
其中 $y'=f'(x')$。此外，
$$
\degtr k(x)/k(y)=\degtr k(x')/k(y'){\quoi},
$$
因为扩张 $k(x')/k(x)$ 和 $k(y')/k(y)$ 都是代数扩张。因此只需证明
$$
d(y')=d(y)-d(y_{0}){\quoi}.
$$

可设 $y$ 是 $Y$ 中的极大点，因而 $y'$ 是 $Y'$ 中的极大点。于是
\begin{align*}
d(y') &= \dim \Oo_{Y',\bar{y}_{0}}\\
      &= \dim \Oo_{Y,y_{0}}\quad \text{\rm (\EGA IV 6.1.3)}\\
      &= \sisi{\dim Y}{\dim_{y_0} Y} -d(y_{0})\quad \text{\rm (\EGA IV 5.2.3.1)}\\
      &= d(y)-d(y_{0}){\quoi},
\end{align*}
结论得证。

\section{仿射代数概形的上同调维数}\etiquette[3]{XIV.sec3}

\begin{enonce*}{定理 3.1~(d)}\etiquette[3.1]{XIV.thm3.1}\nde{定理 \ref{XIV.thm3.1} 现在称为\textit{仿射 Lefschetz 定理}。}
设 $f:X\rightarrow Y$ 是域 $k$ 上有限型概形之间的仿射态射。设 $F$ 是 $X$ 上的挠层，满足 $d(F)\leq d$ \emph{（\cf \hyperref[XIV.subsec2.1]{2.1}）}。那么 $d(R^{q}f_{\ast}F)\leq d-q$。
\end{enonce*}

特别地，取 $Y=\Spec(k)$，得到：

\begin{enonce*}{推论 3.2}\etiquette[3.2]{XIV.cor3.2}
设\pageoriginale $X$ 是可分闭域 $k$ 上的有限型仿射概形。那么\nde{考虑复解析几何中的类似断言：设 $X$ 是 $\mathbf{C}$ 上的仿射簇，$F$ 是 $X$ 上的可构造层。那么对于每个 $m>\dim X$，有 $H^m(X^{\rm an},F)=0$，其中 $X^{\rm an}$ 是与 $X$ 相伴的解析空间。虽然这一断言由 \hyperref[XIV.cor3.2]{3.2} 及比较定理 \hyperrefext[XVI][XVI.thm4.1]{4.1} 得出，但还有两个完全解析的证明，分别由 Pierre Deligne 和 Hélène Esnault 给出；见 Esnault, H., Variation on Artin's vanishing theorem, \textit{Advances in Math.} \textbf{198} (2005), 435--438。}
$$
\cd X\leq \dim X{\quoi}.
$$
\end{enonce*}

\begin{enonce*}{推论 3.3}\etiquette[3.3]{XIV.cor3.3}
设 $X\subset \PP^N_k$ 是可分闭域 $k$ 上维数为 $n$ 的射影概形，$Y=H\cap X$ 是 $X$ 的一个“超平面截面”。那么，对于 $X$ 上的每个挠层 $F$，典范同态
$$
H^{q}_Y(X,F)\longrightarrow H^{q}(X,F)
$$
当 $q>n+1$ 时为双射，当 $q=n+1$ 时为满射。
\end{enonce*}

这一推论直接由 \sisi{\hyperref[XIV.def2.2]{2.2}}{\hyperref[XIV.cor3.2]{3.2}} 及闭子集的正合列（\sisi{IV 3.7.1}{\hyperrefext[V][V.6.5.4]{6.5.4}}）得出，
%%% This last reference should read \ref{V.eq6.5.4}, but the TeX-file of chapter V contains the
%%% label without "eq".
因为 $X-Y=\sisi{\mathcal{U}}{U}$ 是维数 $\leq n$ 的仿射概形。注意，\hyperref[XIV.cor3.3]{3.3} 推广了 Lefschetz 关于超平面截面的一个定理。事实上，若 $X$ 在 $k$ 上光滑，$F$ 局部常值且其阶与特征互素，则利用相对上同调纯性定理（\hyperrefext[XVI][XVI.sec3]{3}）及谱序列（\sisi{\quad}{\hyperrefext[V][V.6.4.3]{6.4.3}}），可以显式计算带支集上同调 $H^{q}_Y(X,F)$，
%%% This last reference should read \ref{V.eq6.4.3}, but the TeX-file of chapter V contains the
%%% label without "eq".
得到
$$
H^{q}_Y(X,\ZZ/n)\simeq H^{q-2}(Y,\mu_n^{-1}) \mbox{\rm\quad 对所有 $q$；}
$$
由 \hyperref[XIV.cor3.3]{3.3} 给出的同态 $H^{q-2}(Y,\mu_n^{-1})\rightarrow H^{q}(X,\ZZ/n)$ 正是“Gysin 同态”，我们将在后面研究它\footnote[1]{关于 Lefschetz 型定理的系统研究，\Cf \SGA 2 XIV。}。

定理 \hyperref[XIV.thm3.1]{3.1}~(d) 等价\nde{更确切地说，正文中的论证给出以下结果：对于 $d'\leq d$ 的断言 \hyperref[XIV.cor3.4]{3.4}~(d') 蕴含 \hyperref[XIV.thm3.1]{3.1}~(d)，而断言 \hyperref[XIV.thm3.1]{3.1}~(d) 蕴含 \hyperref[XIV.cor3.4]{3.4}~(d)。这足以得到下面的结论。}于下面的陈述：

\begin{enonce*}{推论 3.4~(d)}\etiquette[3.4]{XIV.cor3.4}
设 $Y'$ 是严格局部概形，且是某个域上代数概形的严格局部化。设 $f':X'\rightarrow Y'$ 是有限型仿射态射\pageoriginale，$F'$ 是 $X'$ 上满足 $\delta(F')\leq d$ 的挠层 \emph{（\cf （\hyperref[XIV.def2.2]{2.2}））}。那么当 $q>d$ 时，$H^{q}(X',F')=0$。
\end{enonce*}

事实上，设 \hyperref[XIV.cor3.4]{3.4}~(d) 成立。设 $f:X\rightarrow Y$ 是 $k$ 上代数概形之间的仿射态射，$F$ 是 $X$ 上满足 $d(F)\leq d$ 的挠层。设 $y$ 是 $Y$ 的一个点，$\xi$ 是 $y$ 上方的几何点。为证明 \sisi{\ref{XIV.subsec3.1}}{\hyperref[XIV.thm3.1]{3.1}}~(d)，只需证明当 $d(y)>d-q$ 时，茎 $(R^{q}f_{\ast}F)_{\xi}$ 为零。

设 $f':X'\rightarrow Y'$ 是 $f$ 在点 $\xi$ 处的严格局部化；即 $Y'$ 为 $Y$ 在 $\xi$ 处的严格局部化，$X'=\fprod{X}{Y'}{Y}$。

设 $F'$ 是 $F$ 在 $X'$ 上诱导的层，于是 $(R^{q}f_{\ast}F)_{\xi}=H^{q}(X',F')$（\hyperrefext[VIII][VIII.thm5.2]{5.2}）。
\sisi{有
$$
\delta(F')\leq d(F)-d(y)\leq d-d(y)\quad (\hyperref[XIV.prop2.4]{2.4}){\quoi},
$$}
{由（\hyperref[XIV.prop2.4]{2.4}），有
$$
\delta(F')\leq d(F)-d(y)\leq d-d(y){\quoi},
$$}
%%% I find it irritating that the reference to the proposition from which the above follows,
%%% is referred to in a style that might be confused with an equation number
%%% (although the correct equation number is 3.4.1)
并且应用 \hyperref[XIV.cor3.4]{3.4}~(d)，得到当 $q>d-d(y)$ 时 $H^{q}(X',F')=0$，这就证明了结论。

反过来，设 \sisi{\hyperref[XIV.subsec2.1]{2.1}}{\hyperref[XIV.thm3.1]{3.1}}~(d) 成立，$F'$ 和 $f':X'\rightarrow Y'$ 如 \hyperref[XIV.cor3.4]{3.4}~(d) 的陈述所述。由于上同调与滤过正向极限交换（\hyperrefext[VII][VII.prop3.3]{3.3}），且 $F'$ 是其可构造子层的正向极限（\hyperrefext[IX][IX.prop2.9]{2.9}），可设 $F'$ 可构造。现在 $Y'$ 是代数概形 $Y$ 在\emph{闭}点 $y_0$ 上方几何点 $\xi$ 处的严格局部化（\hyperrefext[X][X.lem3.3]{3.3}）。写成 $Y'=\varprojlim_{\alpha}Y_{\alpha}$，其中 $(Y_{\alpha})$ 是 $Y$ 上仿射平展概形的滤过逆向系统（\hyperrefext[VIII][VIII.subsec4.5]{4.5}）。由 \hyperrefext[IX][IX.cor2.7.4]{2.7.4}，当 $\alpha$ 充分大时，$F'$ 由 $\sisi{X=\fprod{X}{Y_K}{Y}}{X_\alpha=\fprod{X}{Y_\alpha}{Y}}$ 上的某个可构造挠层 \sisi{}{$F_\alpha$} 诱导。利用 \hyperref[XIV.prop2.4]{2.4} 和\pageoriginale \EGA IV 9.5.5，立即可见还可设 $d(F_{\alpha})=\delta(F')\leq d$。用 $Y_{\alpha}$ 代替 $Y$，归结为 $F'$ 由 $X$ 上满足 $d(F)\leq d$ 的层 $F$ 诱导的情形。于是，由 \sisi{\hyperref[XIV.subsec2.1]{2.1}}{\hyperref[XIV.thm3.1]{3.1}}~(d)，当 $q>\sisi{D}{d}$ 时 $R^{q}f_{\ast}F=0$，因而特别地
$$
(R^{q}f_{\ast}F)_{\xi}=H^{q}(X',F')=0\quad{\mbox{\rm 若}}\quad
q>d{\quoi},
$$
这就证明了 \hyperref[XIV.cor3.4]{3.4}~(d)。

作为 \hyperref[XIV.cor3.4]{3.4}~(d) 的特例，有：

\begin{enonce*}{推论 3.5}\etiquette[3.5]{XIV.cor3.5}
设 $Y'$ 是维数 $\leq d$ 的严格局部概形，且是域 $k$ 上某个代数概形的严格局部化。设 $U\subset Y'$ 为仿射开子集。那么 $\cd U\leq d$。
\end{enonce*}

\begin{enonce*}[remark]{注 3.6}\etiquette[3.6]{XIV.rem3.6}
很有可能\nde{Ofer Gabber 已成功证明 \ref{XIV.thm3.1} 的以下类似结果：设 $f:X\to Y$ 是有限型仿射态射，其中 $Y$ 拟优（\EGA IV 7.8），并赋有维数函数 $\delta_Y$。设 $\delta_X$ 为 $X$ 上相应的修正维数函数。若 $n$ 在 $Y$ 上可逆，那么对于 $X$ 上的每个 $\mathbf{Z}/n$-模层 $F$，以及每个 $q\geq0$，有 $\delta_Y(R^q f_\ast F)\leq\delta_X(F)-q$。这里与前面一样，对层 $G$ 定义 $\delta(G)=\sup\{\delta(x)\mid x\in X\quaet G_{\bar{x}}\neq0\}$。
}，\hyperref[XIV.cor3.4]{3.4}，从而 \hyperref[XIV.cor3.5]{3.5}，对于\emph{每个} Noether 局部概形 $Y'$（不必是代数概形的严格局部化）仍然成立，至少当 $Y'$ 是\emph{优的}（\EGA IV 7.8）时如此。在进一步假设 $Y$ 的特征为零时，将在（\sisi{XX 6}{\hyperrefext[XIX][XIX.sec6]{6}}）中证明这一点。
\end{enonce*}

\section[定理 3.1 的证明]{定理 \texorpdfstring{\hyperref[XIV.thm3.1]{3.1}}{3.1} 的证明}\etiquette[4]{XIV.sec4}

证明对 $d$ 作归纳。

\begin{enonce*}{引理 4.1}\etiquette[4.1]{XIV.lem4.1}
\emph{\hyperref[XIV.thm3.1]{3.1}~(0)} 和 \emph{\hyperref[XIV.thm3.1]{3.1}~(1)} 成立。
\end{enonce*}

\noindent
\emph{证明}。当 $d=0$ 时这是显然的，此时 $\supp F$ 在 $k$ 上有限，应用（\hyperrefext[VIII][VIII.prop5.5]{5.5}）即可。设 $F$ 和 $f:X\rightarrow Y$ 如 \hyperref[XIV.thm3.1]{3.1} 的陈述所述，并设 $d(F)\leq1$。用 $F$ 的一个可构造子层代替它（\hyperrefext[IX][IX.prop2.9]{2.9} 和 \hyperrefext[VII][VII.prop3.3]{3.3}），便归结为 $F$ 可构造的情形。于是 $F$ 的支集\pageoriginale 包含于 $X$ 的一个维数 $\leq1$ 的闭子概形中（\hyperrefext[IX][IX.def2.3]{2.3}\sisi{ (ii)}{}），可以用这个子概形代替 $X$；即可以假设 $\dim X\leq1$。由 \hyperrefext[IX][IX.cor5.7]{5.7}，可分闭域上维数为 $1$ 的仿射代数概形的上同调维数 $\leq1$。此外，只需处理 $k$ 可分闭的情形。事实上，设 $\bar f:\bar X\rightarrow\bar Y$ 为由 $f$ 经基变换 $\Spec\bar k\rightarrow\Spec k$ 得到的态射，其中 $\bar k$ 是 $k$ 的可分闭包。对于 $Y$ 的几何点 $\xi$，$f$ 在 $\xi$ 处的严格局部化 $f':X'\rightarrow Y'$ 可视为 $\bar f$ 的一个严格局部化，因而纤维 $(R^{q}f_{\ast}F)_\xi$ 可视为 $R^{q}\bar f_{\ast}\bar F$ 的一个纤维。由于 $f$ 仿射，得到当 $q>1$ 时 $R^{q}f_{\ast}F=0$。显然还有 $d(f_{\ast}F)\leq1$。对于 $q=1$，注意假设 $\dim X\leq1$ 蕴含 $Y$ 有一个有限闭子集 $Z$，使 $f$ 在 $Y-Z$ 上有限。由此，$R^{1}f_{\ast}F$ 在 $Y-Z$ 上为零（\hyperrefext[VIII][VIII.cor5.6]{5.6}），因而 $d(R^{1}f_{\ast}F)\leq1$。引理得证。

设 $d\geq2$，现在假设对于 $d'<d$，\hyperref[XIV.thm3.1]{3.1}~(d')（或等价地，\hyperref[XIV.cor3.4]{3.4}~(d')）成立。

\begin{enonce*}{引理 4.2}\etiquette[4.2]{XIV.lem4.2}
为在一般情形证明 \emph{\hyperref[XIV.cor3.4]{3.4}~(d)}，只需处理 $X'=\EE^1_{Y'}$ 是 $Y'$ 上的仿射直线，且 $f$ 是结构态射的情形。
\end{enonce*}

\noindent
\emph{证明}。显然可设对于某个适当的 $N$，有 $X'=\EE^N_{Y'}$，因为 $X'$ 可嵌入某个 $\EE^N_{Y'}$。设 $N>1$，且对于每个 $Y'$ 和 $r<N$，结论对于 $\EE^r_{Y'}$ 已知。设 $F'$ 是 $\EE^N_{Y'}$ 上的挠层，满足\pageoriginale $d(F')\leq d$。考虑图
$$
\xymatrix@=1.5cm{
\EE^{N}_{Y'} \ar[rd]_{f'}       \ar[rr]^{g'}&      &
\EE^{N-1}_{Y'} \ar[ld]^{h'}\\
                       & Y' & &,
}
$$
其中 $g'$ 和 $h'$ 为典范投影，所以 $\EE^N_{Y'}$ 同构于 $\EE^{N-1}_{Y'}$ 上的仿射直线。由于 $\EE^{N-1}_{Y'}$ 的每个严格局部化显然也是某个代数概形的严格局部化，可对 $R^{q}g'_{\ast}F'$ 的纤维应用归纳假设，得到
\begin{equation*}
\delta(R^{q}g'_{\ast}F')\leq d-q{\quoi}.\tag{$\ast$}
\end{equation*}
此外，由关于 $N$ 的归纳假设，\hyperref[XIV.cor3.4]{3.4}~(d') 对 $h'$ 和 $d'\leq d$ 成立。

考虑 Leray 谱序列
$$
E^{p,q}_2=H^{p}(\EE^{N-1}_{Y'},R^{q}g'_{\ast}(F'))
\Longrightarrow H^{n}(\EE^N_{Y'},F'){\quoi}.
$$
对态射 $h'$ 应用 $(\ast)$ 和 \hyperref[XIV.cor3.4]{3.4}~(d')，得到当 $p>d-q$ 时 $E^{p,q}_2=0$。因此收敛项在 $n>d$ 时为零，引理得证。

考虑下面的附加陈述，它是 \hyperref[XIV.cor3.5]{3.5} 的一个特例：

\begin{enonce*}{陈述 4.3~(d)}\etiquette[4.3]{XIV.enon4.3}
设 $X'$ 严格局部，维数 $\leq d$，并且是某个域上代数概形的严格局部化。设 $f\in\Gamma(X',\Oo_{X'})$，$U\subset X'$ 为开子集 $U=X'-V(f)$。那么 $\cd U\leq d$。
\end{enonce*}

\begin{enonce*}{引理 4.4}\etiquette[4.4]{XIV.lem4.4}
\emph{\hyperref[XIV.enon4.3]{4.3}~(d)} 蕴含 \emph{\hyperref[XIV.cor3.4]{3.4}~(d)}。
\end{enonce*}

\noindent
\emph{证明}。\pageoriginale 考虑图
$$
\xymatrix@=1.5cm{
\EE^1_{Y'} \ar[rr]^i \ar[rd]_{f'}&            &
\PP^1_{Y'} \ar[ld]^{g'}\\
                    &       Y'   &
}
$$
其中 $Y'$ 如 \sisi{2.6}{\hyperref[XIV.cor3.4]{3.4}} 的陈述所述，而 $\PP^1_{Y'}$ 是由 $\EE^1$“添上无穷远截面 $Y'_\infty$”得到的射影直线。设 $F'$ 是 $\EE^1$ 上满足 $d(\supp F')\leq d$ 的挠层。有谱序列
$$
H^{p}(\PP^1,R^{q}i_{\ast}F')\Longrightarrow
H^{n}(\EE^1,F'){\quoi}.
$$

当 $q>0$ 时，$R^{q}i_{\ast}F'$ 的支集在 $Y'_\infty$ 中，而它与 $Y'$ 为 $Y'$-同构。由于 $Y'$ 严格局部，得到当 $p,q>0$ 时 $H^{p}(\PP^1,R^{q}i_{\ast}F')=0$。此外，由\emph{固有基变换定理}（\hyperrefext[XII][XII.thm5.1]{5.1}），$\PP^1$ 上挠层的上同调可在闭纤维上计算，而闭纤维是上同调维数 $\leq2$ 的概形（\hyperrefext[X][X.cor4.3]{4.3}）。因此当 $p>2$ 时 $H^{p}(\PP^1,i_{\ast}F')=0$。由于我们要证明当 $n>d$ 时 $H^{n}(\EE^1,F')=0$，且 $d\geq2$，便归结为证明当 $q>d$ 时 $H^{0}(\PP^1,R^{q}i_{\ast}F')=0$。这一群还同构于 $H^{0}(Y'_\infty,R^{\sisi{1}{q}}i_{\ast}F')$，由（\hyperrefext[VIII][VIII.prop4.6]{4.6}），它同构于 $R^{q}i_{\ast}F'$ 在点 $Q$ 处的纤维，其中 $Q$ 是 $\PP^1$ 在闭纤维中的无穷远点。

只需处理 $F'$ 可构造的情形（\hyperrefext[VII][VII.prop3.3]{3.3} 和 \hyperrefext[IX][IX.prop2.9]{2.9}）。由于 $\delta(\supp F')\leq d$，于是存在 $\PP^1$ 的一个闭子概形 $X'$，满足 $\dim X'\leq d$，并包含 $F'$ 的支集。设 $\widetilde X'$ 为 $X'$ 在几何点 $Q$ 处的局部化，令
$$
\widetilde U=\widetilde X'-
\fprod{\widetilde X'}{Y'_\infty}{\PP^1},
$$
并设 $\widetilde F$ 为 $F'$ 经态射 $\widetilde U\rightarrow\PP^1$ 在 $\widetilde U$ 上诱导的层\pageoriginale。那么 $R^{q}i_{\ast}F'$ 在 $Q$ 处的纤维恰为 $H^{\sisi{p}{q}}(\widetilde U,\widetilde F)$（\hyperrefext[VIII][VIII.thm5.2]{5.2}），对严格局部概形 $\widetilde X'$ 及开子集 $\widetilde U$ 应用 \hyperref[XIV.enon4.3]{4.3}~(d) 即可。

以下引理与 \hyperref[XIV.lem4.4]{4.4} 一起将完成证明。

\begin{enonce*}{引理 4.5}\etiquette[4.5]{XIV.lem4.5}
对于 $d'<d$ 的 \emph{\hyperref[XIV.cor3.4]{3.4}~(d')} 蕴含对于 $d'\leq d$ 的 \emph{\hyperref[XIV.enon4.3]{4.3}~(d')}。
\end{enonce*}

\noindent
\emph{证明}。设 \hyperref[XIV.cor3.4]{3.4}~(d') 对 $d'<d$ 成立。由归纳法，可设 \hyperref[XIV.enon4.3]{4.3}~(d') 对 $d'<d$ 成立。设 $X'$、$f$ 和 $U$ 如陈述 \hyperref[XIV.enon4.3]{4.3}~(d) 所述。于是 $X'$ 是域 $k$ 上代数概形 $X_0$ 在一个\emph{闭}几何点处的严格局部化。用 $x_0$ 的一个邻域代替 $X_0$ 之后，可取 $X_0$ 仿射，维数 $\leq d$。于是 $X'$ 是仿射概形 $X_i$ 的逆向极限，这些概形在 $X_0$ 上平展（\hyperrefext[VIII][VIII.subsec4.5]{4.5}）。

设 $F$ 是 $U$ 上的挠层。须证当 $q>d$ 时 $H^q(U,F)=0$，而仍然只需处理 $F$ 可构造的情形。于是资料 $f$、$F$ 和 $U$ 来自某个 $X_i$ 上的类似资料（\sisi{\hyperrefext[IX][IX.prop2.7]{2.7}}{\hyperrefext[IX][IX.cor2.7.4]{2.7.4}}）。因此存在 $U_i=X_i-V(f_i)$ 上的层 $F_i$，它在 $U$ 上诱导 $F$。由于 $U=\varprojlim U_i$，有（\hyperrefext[VII][VII.cor5.8]{5.8}）
$$
H^q(U,F)=\varinjlim H^q(U_i,F_i){\quoi}.
$$
所以只需证明以下事实：对于每个 $i$，存在一个交换图
$$
\xymatrix@=1.5cm{
 U \ar[rr] \ar[rd]&        & U''\ar[ld]\\
   &  U_i   &
}
$$
其中 $U\rightarrow U_i$ 为给定的箭头，且 $\cd(U'')\sisi{<}{\leq}d$。事实上，这便蕴含\pageoriginale 当 $q>d$ 时，$H^q(U_i,F_i)$ 的每个元素在 $H^q(U,F)$ 中的像都为零。

设 $i=0$，考虑由 $f_0\in\Gamma(X_0,\Oo_{X_0})$ 定义的态射 $g:X_0\rightarrow\EE^1_k=Y_0$。设 $Y''$ 为 $\EE^1_k$ 在原点上方某个几何点处的严格局部化。于是 $Y''$ 为离散赋值环的谱。令 $X''=\fprod{X_0}{Y''}{Y_0}$。

由于 $X'$ 严格局部，态射 $X'\rightarrow X_0$ 通过 $X''$ 分解。此外，开子集 $U''=\fprod{X''}{U_0}{X_0}$ 恰为 $X''/Y''$ 的泛纤维。因此，它是 $Y''$ 的泛点上的仿射代数概形，而且显然 $\dim U''\leq d-1$。

验证 $\cd U''\leq d$，就完成了证明。设 $K$ 为 $Y''$ 在其泛点处的剩余域。这个域的上同调维数为 $1$；即 $\cd(G)=1$，其中 $G$ 为 $K$ 的可分闭包 $\bar K$ 的 Galois 群（\hyperrefext[X][X.thm2.2]{2.2}）。令 $\bar U''=\fprod{U''}{(\Spec\bar K)}{\Spec K}$。由关于 $d$ 的归纳假设，有 $\cd\bar U''\leq d-1$（这是\sisi{对维数 $\leq d$ 的概形应用 \hyperref[XIV.def2.2]{2.2}}{对维数 $<d$ 的概形应用 \hyperref[XIV.cor3.2]{3.2}}）。于是，由 Hochschild--Serre 谱序列（\hyperrefext[VIII][VIII.prop8.4]{8.4}），得 $\cd U''\leq d$：
$$
H^p(G,H^q(\bar U\sisi{'}{''},\bar F))\Longrightarrow H^n(U'',F){\quoi},
$$
这就证明了引理 \hyperref[XIV.lem4.5]{4.5}。
